\section{Superlinear ordinary correlated-agreement failures in quadratic extensions}
\label{app:ordinary-ca-superlinear}

The quadratic lower bound of Corollary~\ref{cor:quadratic-extension-fixed-gap}
concerns full agreement-set MCA. A separate construction excludes even
ordinary correlated agreement, at the cost of a smaller exponent and a
gap that is bounded below rather than fixed exactly. We make no novelty
claim for the standard sieve input or the combination below.

\begin{theorem}\label{thm:ordinary-ca-superlinear}
For every $0<\varepsilon<1/8$, there are unbounded lengths $n$ and primes
$p=12n/13+5$ for which a Reed--Solomon code over $\F_{p^2}$, of exact rate
$3/13$, has a received affine line with
$\Omega_\varepsilon(n^{5/4-\varepsilon})$ nearby labels at an agreement
threshold $M$, but no ordinary correlated agreement of size $M$.
The capacity gap satisfies $M/n-3/13>3/26$, and the radius is strictly
below the characteristic-based Elias radius for all sufficiently large
members of the family.
\end{theorem}

Thus a linear ordinary-CA bound uniform over gaps bounded below is
impossible in this large-characteristic field class. The actual ratio
$M/n$ is not asserted to be constant. This is neither a prime-ambient-field
result nor a first-order-regime tightness statement.

\paragraph{The prime-selection input.}
Put $\theta=1/4-\varepsilon$. The lower linear sieve and
Bombieri--Vinogradov imply that there are infinitely many primes
$p\equiv41\pmod{48}$ for which every odd prime divisor of $(p-1)/4$
is at least $p^\theta$. Here is the specific reduction, so that the fixed
progression introduces no hidden distribution hypothesis. Sift
\[
 \mathcal A=\{(p-1)/8:x/2<p\le x,\ p\equiv41\pmod{48}\}.
\]
For squarefree $d$ coprime to $6$, the condition $d\mid(p-1)/8$
selects one reduced residue class modulo $48d$. Hence the local density
is $g(d)=1/\varphi(d)$, with main size
$X=(\operatorname{Li}(x)-\operatorname{Li}(x/2))/16$; set $g(2)=g(3)=0$.
Bombieri--Vinogradov bounds the sum of absolute remainders through
$D=x^{1/2-\varepsilon}$ by $O_A(x/\log^A x)$ for every fixed $A$.
With $z=x^{1/4-\varepsilon}$, the sieve parameter
\[
 s=\frac{\log D}{\log z}
   =\frac{1/2-\varepsilon}{1/4-\varepsilon}\in(2,3)
\]
has positive lower sieve function
$f(s)=2e^\gamma\log(s-1)/s$. Since
$\prod_{\ell<z}(1-g(\ell))\asymp1/\log z$, there are
$\gg_\varepsilon x/\log^2x$ unsifted primes. Their associated odd
integers $(p-1)/8$ have no prime factor below $z\ge p^\theta$.
This is the standard shifted-prime sieve argument; see
\cite[Section~5, Example~3, pp.~36--37]{heath-brown-sieves} and
\cite[Theorem~2 and equation~(11)]{tao-linear-sieve}.

\paragraph{Excluding small polynomial orbits.}
Let $\zeta$ be a primitive eighth root of unity. No affine polynomial
agrees at three distinct points of $\mu_8$ with
\[
 W_2(Y)=(1+Y^4)/2-Y^2
\]
in characteristic $p>17$ with $p\equiv1\pmod8$. For completeness this
finite assertion can be checked integrally. Write
$v_j=1+\zeta^{4j}-2\zeta^{2j}$. For $0\le i<j<h<8$, take the norm of
\[
 (\zeta^j-\zeta^i)(v_h-v_i)
 -(\zeta^h-\zeta^i)(v_j-v_i)
\]
from $\mathbb Q(\zeta)$ to $\mathbb Q$, equivalently its resultant
with $Z^4+1$ after reducing powers of $Z$. The complete distribution is
\[
\begin{array}{c|rrrrrrr}
 |\mathrm{Norm}|&64&256&576&1088&3136&4096&6400\\\hline
 \text{number of triples}&12&8&8&8&4&8&8.
\end{array}
\]
All $56$ determinants are nonzero and their norms have prime factors only
$2,3,5,7,17$. Reduction at $p>17$ therefore preserves nonvanishing.
An independent exact verifier and the individual determinants are in
\path{research/dickson_fixed_gap/verify_two_orbit.py} and its JSON output.

\begin{proof}[Proof of Theorem~\ref{thm:ordinary-ca-superlinear}]
Choose the primes above, put $k=(p-1)/4$ and $N=4k$, and use
$W(X)=(1+X^{2k})/2-X^k$ on $\F_p^\times$. Let $M$ be its true maximum
agreement with degree-$<k$ polynomials. Proposition~\ref{prop:full-length-fixed-gap}
gives $M\ge3k/2$, while the degree of $W-P$ gives $M\le2k$.
The group $H=\mu_k$ preserves the word and acts on its nearest list
by $P(X)\mapsto P(hX)$.

The size $r$ of any such orbit divides $k$. Here $v_2(k)=1$, so the
roughness condition implies either $r\in\{1,2\}$ or $r\ge p^\theta$.
In the former case invariance under a subgroup of order $k/2$ forces
$P(X)=a+bX^{k/2}$. Passing to $Y=X^{k/2}$, whose fibers have size $k/2$,
the preceding eight-point calculation bounds its agreement by $k$.
Such a polynomial cannot be nearest. Consequently there are at least
$p^\theta$ nearest polynomials. Select $L=\lfloor p^\theta\rfloor$
of them; for large $p$, $L\le k$.

This nearest agreement remains $M$ over $\F_{p^2}$: a degree-$<k$
polynomial matching $k$ base-field coordinates and values has base-field
coefficients by interpolation. Choose an agreement anchor incident to
$\ell\ge ML/N$ selected candidates, and divide each candidate minus the
received anchor value by the anchor's linear factor. On the $N-1$ old
coordinates the resulting degree-$\le k-2$ quotients each have exactly
$M-1$ agreements. No polynomial of this degree has $M$ old agreements,
since restoring the anchor would contradict maximality of $M$.

Let $Q=p^2$ and $q=(k-10)/3$. This is an integer because
$p\equiv41\pmod{48}$ gives $k\equiv1\pmod3$. For each of the $Q-N$
unused field points let $s_x$ count distinct quotient values. Pairwise
root counting and Cauchy--Schwarz give average diversity at least
\[
 \mu=\frac{\ell(Q-N)}{Q-N+(\ell-1)(k-2)}\ge\ell/2,
\]
where $Q\ge2NL$. Choose the $q$ points with largest $s_x$, give them
direction one, and retain direction zero on the old domain. Independent
uniform translations of the new value sets have expected union size
\[
 Q\left[1-\prod_x(1-s_x/Q)\right]
 \ge Q[1-(1-\mu/Q)^q]\ge q\ell/4\ge\frac{qML}{4N}.
\]
For the middle inequality, use $1-u\le e^{-u}$ and
$q\ell/(2Q)\le1/4$. Every counted label has at least $M$ agreements
with a selected quotient.

If a correlated witness pair $F,G$ of degree at most $k-2$ has $G=0$,
its joint agreements are confined to the old domain, with maximum
$M-1$. If $G\ne0$, there are at most $k-2$ old zeros and at most $q$
new agreements. Since $q\le M-k+1$, this again totals at most $M-1$.
Thus ordinary correlated agreement of size $M$ is absent.

Finally,
\[
 K=k-1,\qquad n=N-1+q=\frac{13(k-1)}3,\qquad
 \frac{M-K}{n}\ge\frac{k/2+1}{n}>\frac3{26}.
\]
For $k\ge20$ we have $q\ge k/6$, so the number of nearby labels is at least
\[
 \frac{qML}{4N}\ge\frac{kL}{64}
 \ge\frac{3nL}{832}\ge\frac{nL}{300}
 =\Omega_\varepsilon(n^{5/4-\varepsilon}).
\]
Here $p=12n/13+5>K-1$. Also
$H_p(1-M/n)\le1-M/n+1/\log_2p<1-K/n$ for large $p$,
which proves the strict Elias assertion.
\end{proof}

A complete small-field replay in
\path{research/ordinary_ca_superlinear/verify.py} exhausts all $17^4$
source polynomials and all $17^3$ anchored-code polynomials. It certifies
$27$ nearby labels at threshold $6$ for length $18$ and dimension $3$
over $\F_{17^2}$, with old and nonzero-direction joint-agreement bounds
both equal to $5$. This checks the boundary mechanism, not the
asymptotic sieve theorem or the Elias condition for this small fixture.
